\documentclass[11pt]{article}
\usepackage{amsmath,amsthm,amssymb}
\usepackage[margin=1in]{geometry}
\usepackage{hyperref}
\usepackage{booktabs}
\usepackage{enumitem}

\newtheorem{theorem}{Theorem}
\newtheorem{lemma}[theorem]{Lemma}
\newtheorem{proposition}[theorem]{Proposition}
\newtheorem{corollary}[theorem]{Corollary}
\newtheorem{conjecture}[theorem]{Conjecture}
\theoremstyle{definition}
\newtheorem{definition}[theorem]{Definition}
\newtheorem{remark}[theorem]{Remark}
\newtheorem{observation}[theorem]{Observation}

\newcommand{\Hq}{H_q}
\newcommand{\Sn}{S_n}
\newcommand{\Vlam}{V_\lambda}
\newcommand{\Vmu}{V_\mu}
\newcommand{\Bdec}{B^{(\lambda)}}
\newcommand{\Om}{\Omega}
\newcommand{\hatl}{\widehat\lambda}
\newcommand{\hatm}{\widehat\mu}
\newcommand{\hatn}{\widehat\nu}
\newcommand{\supp}{\mathrm{supp}}
\newcommand{\SYT}{\mathrm{SYT}}
\newcommand{\Rp}{{R'}}
\newcommand{\elh}{\ell(\hatl)}
\newcommand{\Outer}{\mathcal{O}}

\title{Sharpness of the multi-row containment at $\tau+1$:\\
       \large $B^{(\lambda)} \cap E^-_{\tau+1}\!\mid_{\Vlam} = 0$}
\author{Clio}
\date{17 May 2026}

\begin{document}
\maketitle

\begin{abstract}
Let $\hatl$ be a partition with $\elh \ge 2$ and every row of length $\ge 2$,
and let $\lambda = (\hatl, 1^q) \vdash n$. The Pillar~1 meta-theorem
(15 May 2026) proves that the column-$1$ chain image
$\Bdec = \Om^{(\lambda)} \Vlam$ is contained in
$\bigcap_{j=1}^{\tau} E_j^-|_{\Vlam}$, where
$\tau = \tau(\hatl) := \max(0,\, q + 2\elh - 1 - |\hatl|)$, for
$q \ge |\hatl| - 2\elh + 2$. The structural sharpness of this
containment was left open. We prove the strong sharpness
\[
\Bdec \cap E^-_{\tau+1}\!\mid_{\Vlam} \;=\; \{0\},
\]
equivalently, the operator $S_{\tau+1} := T_{\tau+1} + 1$ is
\emph{injective} on $\Bdec$. The proof is a strong induction on
$|\hatl|$ that lifts cleanly through Pillar~1's reduction theorem, with
two ingredients: (i)~the outer product
$\Outer = \prod_{j=\ell}^{n-2}(T_j+1)$ is \emph{injective} on
$E^+_{n-1}|_{\Vlam}$ by adjacent-injectivity of the
$+q$-eigenspace chain (May-19 Phase~A), and (ii)~$S_{\tau+1}$ commutes
with every factor of $\Outer$ and with every Phi-embedding, by the
geometric inequality $|j - (\tau+1)| \ge 2$ that holds throughout the
standing regime. The hook base layer ($\hatl$ of length $1$) is handled
by a separate induction on row length $k_h$ with an explicit
$2$-block computation at $k_h = 2$.
\end{abstract}

\section{Setup and statement}\label{sec:setup}

We adopt the conventions of the Pillar~1 paper~\cite{pillar1}.
$\Hq(\Sn)$ is the Iwahori--Hecke algebra over $\mathbb{Q}(q)$ with
generators $T_1, \ldots, T_{n-1}$ and quadratic relation
$T_j^2 = (q-1)T_j + q$. $\Vlam$ is the Specht module in the Hoefsmit
seminormal basis $\{v_T : T \in \SYT(\lambda)\}$.
Set $S_j := T_j + 1$ for $1 \le j \le n-1$, and write
\[
   E_j^+ \;:=\; \ker(T_j - q)\!\mid_{\Vlam}
   \;=\; \mathrm{Im}(S_j\!\mid_{\Vlam}),
   \qquad
   E_j^- \;:=\; \ker(S_j)\!\mid_{\Vlam}.
\]
For $k \ge 1$, $\Rp_k := S_{k-1} S_{k-2} \cdots S_1$, with
$\Rp_1 = 1$. For $\nu \vdash N$ with $\ell = \ell(\nu)$,
\[
   \Om^{(\nu)} \;:=\; \Rp_{\ell+1}^{(\nu)} \Rp_{\ell+2}^{(\nu)} \cdots
   \Rp_N^{(\nu)}, \qquad B^{(\nu)} \;:=\; \Om^{(\nu)} V_\nu.
\]

Throughout, $\hatl$ denotes a partition with $\elh \ge 2$ rows, every
row $\hatl_r \ge 2$, and $\lambda = (\hatl, 1^q) \vdash n$, $n = q +
|\hatl|$, $\ell := \ell(\lambda) = \elh + q$. The unfloored threshold
is
\[
   \tau_0(\hatl) \;:=\; q + 2\elh - 1 - |\hatl|,
   \qquad
   \tau(\hatl) \;:=\; \max\bigl(0,\,\tau_0(\hatl)\bigr).
\]

\paragraph{Standing hypothesis.} $q \ge |\hatl| - 2\elh + 2$
(equivalently $\tau_0 \ge 1$).

\begin{theorem}[Strong sharpness]\label{thm:main}
For $\lambda = (\hatl, 1^q)$ with $\hatl$ as above under the standing
hypothesis,
\[
   B^{(\lambda)} \,\cap\, E^-_{\tau+1}\!\mid_{\Vlam} \;=\; \{0\},
\]
equivalently, $S_{\tau+1}\!\mid_{B^{(\lambda)}}$ is injective.
\end{theorem}

\begin{corollary}[Sharpness]\label{cor:sharp}
$B^{(\lambda)} \not\subseteq E^-_{\tau+1}\!\mid_{\Vlam}$. Combined
with Pillar~1, this gives
$\sup\{j : B^{(\lambda)} \subseteq E^-_j\!\mid_{\Vlam}\} = \tau(\hatl)$
exactly.
\end{corollary}

\section{Inputs}\label{sec:inputs}

We collect three structural inputs.

\begin{proposition}[Adjacent-injectivity chain]\label{prop:adj-inj}
For every $\Hq$-module $V$ and every $1 \le j \le |\mathrm{ind}|-2$,
$(T_j+1)\!\mid_{E^+_{j+1}(V)}$ is injective with image in $E^+_j(V)$.
\end{proposition}

\begin{proof}
Image lands in $E^+_j$: by $T_j(T_j+1) = q(T_j+1)$, so
$T_j(T_j+1)v = q(T_j+1)v$, i.e., $(T_j+1)v \in \ker(T_j - q) = E^+_j$.
Injectivity: this is the May-19 Phase~A adjacent-disjointness statement
$E_j^-(V) \cap E_{j+1}^+(V) = 0$, holding for any
$\Hq$-module by the $H_q(S_3)$-irreducible obstruction. So if
$v \in E^+_{j+1}$ and $(T_j+1)v = 0$, then $v \in E^-_j \cap E^+_{j+1}
= 0$.\qed
\end{proof}

\begin{lemma}[Outer product is injective on $E^+_{n-1}$]\label{lem:Outer-inj}
Set $\Outer := \prod_{j=\ell}^{n-2}(T_j+1)$ (product in increasing-$j$
order from left to right). Then $\Outer\!\mid_{E^+_{n-1}|_{\Vlam}}$ is
injective, with image in $E^+_\ell|_{\Vlam}$.
\end{lemma}

\begin{proof}
Read the product as $\Outer\,v = (T_\ell+1)\cdots(T_{n-2}+1)\,v$,
applied right-to-left. For $v \in E^+_{n-1}$,
$(T_{n-2}+1)v \in E^+_{n-2}$ by Proposition~\ref{prop:adj-inj}, and the
restriction $(T_{n-2}+1)\!\mid_{E^+_{n-1}}$ is injective. Iterating: at
each step $j = n-3, n-4, \ldots, \ell$, the input lies in $E^+_{j+1}$
(by the prior step), $(T_j+1)$ maps injectively into $E^+_j$. After
all factors, the image lies in $E^+_\ell$ and the composite is
injective.\qed
\end{proof}

\begin{theorem}[Pillar~1 reduction theorem~\cite{pillar1}]\label{thm:pillar1-red}
Under the standing hypothesis,
\[
   \Bdec \;=\; \Outer \,\cdot\, \sum_{k \in K_{\hatl}}
   \Phi_{k,*}\bigl(B^{(\mu_{k,*})}\bigr),
\]
where $\Phi_{k,*} : V_{\mu_{k,*}} \hookrightarrow \Vlam$ is the
contributor Phi-embedding (Pillar~1, Lemma~4), $T_i$-equivariant for
$1 \le i \le n-3$, and the Phi-images have pairwise disjoint
seminormal supports.
\end{theorem}

\begin{theorem}[Pillar~1 r-additivity~{\cite[Thm.~A]{r-add}}]\label{thm:r-add}
$\dim \Bdec = \sum_{k \in K_{\hatl}} \dim B^{(\mu_{k,*})}$.
Equivalently, $\Outer \circ \bigoplus_k \Phi_{k,*}$ is injective on
$\bigoplus_k B^{(\mu_{k,*})}$.
\end{theorem}

\noindent
Theorem~\ref{thm:r-add} is in fact strengthened by
Lemma~\ref{lem:Outer-inj}: $\Outer \circ \bigoplus_k \Phi_{k,*}$ is
injective on the larger space $\bigoplus_k V_{\mu_{k,*}}$, since the
Phi-images sit in $E^+_{n-1}$ with disjoint supports, and $\Outer$ is
injective there.

\section{Commutativity geometry}\label{sec:commutativity}

The proof rests on two commutativity facts. Let
\[
   \delta \;:=\; \ell - (\tau+1) \;=\; |\hatl| - \elh.
\]

\begin{lemma}[Geometric inequalities]\label{lem:geometric}
Under the standing hypothesis,
\begin{enumerate}[label=(\roman*),topsep=0.2em,itemsep=0.1em]
\item $\delta = |\hatl| - \elh \;\ge\; \elh \;\ge\; 2$.
\item $n - 1 - (\tau+1) = 2|\hatl| - 2\elh - 1 \;\ge\; 2\elh - 1 \;\ge\; 3$.
\item $\tau + 1 \le \ell - 2$ and $\tau + 1 \le n - 4$.
\end{enumerate}
\end{lemma}

\begin{proof}
(i) Since every row of $\hatl$ has length $\ge 2$,
$|\hatl| \ge 2\elh$, so $\delta = |\hatl| - \elh \ge \elh$. By
hypothesis $\elh \ge 2$. (ii) $n - 1 - (\tau+1) = (q + |\hatl| - 1) -
(q + 2\elh - |\hatl|) = 2|\hatl| - 2\elh - 1 \ge 2 \cdot 2\elh -
2\elh - 1 = 2\elh - 1 \ge 3$. (iii) $\tau + 1 = \ell - \delta \le
\ell - 2$, and $\tau + 1 = n - 1 - (2|\hatl| - 2\elh - 1) \le n - 4$
from (ii).\qed
\end{proof}

\begin{corollary}[Commutativity]\label{cor:commutativity}
Under the standing hypothesis,
\begin{enumerate}[label=(\alph*),topsep=0.2em,itemsep=0.1em]
\item $S_{\tau+1}$ commutes with every factor of $\Outer$ (since each
   factor index $j \ge \ell \ge \tau + 3$, so $|j - (\tau+1)| \ge 2$).
\item $S_{\tau+1}$ commutes with $T_{n-1}$ (since $|n - 1 - (\tau+1)|
   \ge 3 \ge 2$), hence preserves $E^+_{n-1}|_{\Vlam}$.
\item $S_{\tau+1}$ commutes with each $\Phi_{k,*}$
   (since $\tau + 1 \le n - 3$, the $T_i$-equivariance range of
   $\Phi_{k,*}$).
\end{enumerate}
\end{corollary}

\section{Proof of the main theorem}\label{sec:proof}

\begin{proof}[Proof of Theorem~\ref{thm:main}]
Strong induction on $|\hatl|$.

\medskip\textbf{Base case ($|\hatl| = 4$, $\hatl = (2, 2)$).}
The unique element of $K_{\hatl}$ is $k = 2$ (shrinking corner), and
the inner contributor shape is the hook $\mu_{2,*} = (2, 1^q)$, with
$\hat\mu = (2)$, $\ell(\hat\mu) = 1$ (outside the standing hypothesis).
By Lemma~\ref{lem:hook-sharp} below, $S_{q}$ is injective on
$B^{(2, 1^q)}$, with $\tau(\hat\mu) + 1 = q = \tau(\hatl) + 1$
(threshold-matching at the hook layer; see
Remark~\ref{rem:hook-tau-match}). The induction step below applies
verbatim with $K_{\hatl} = \{2\}$ and IH replaced by
Lemma~\ref{lem:hook-sharp}.

\medskip\textbf{Inductive step ($|\hatl| \ge 5$).}
Suppose Theorem~\ref{thm:main} holds for every $\hatm$ with
$\ell(\hatm) \ge 2$, every row $\ge 2$, and $|\hatm| < |\hatl|$, and
suppose Lemma~\ref{lem:hook-sharp} holds for every hook $(k_h, 1^{m_h})$
with $k_h \ge 2$, $m_h \ge k_h$. Both inputs are available, the
former by IH and the latter by Lemma~\ref{lem:hook-sharp} (proved
independently in Section~\ref{sec:hook}).

Let $w \in \Bdec$ with $S_{\tau+1} w = 0$. We must show $w = 0$.

By the reduction theorem (Theorem~\ref{thm:pillar1-red}),
\[
   w \;=\; \Outer \,\Bigl(\sum_{k \in K_{\hatl}}
   \Phi_{k,*}(u_k)\Bigr)
   \quad \text{for some } u_k \in B^{(\mu_{k,*})}.
\]
The representation is unique by Theorem~\ref{thm:r-add}.

By Corollary~\ref{cor:commutativity}, $S_{\tau+1}$ commutes with
$\Outer$ and with each $\Phi_{k,*}$. So
\[
   0 \;=\; S_{\tau+1} w
   \;=\; \Outer\,\Bigl(\sum_{k \in K_{\hatl}}
   \Phi_{k,*}(S_{\tau+1} u_k)\Bigr).
\]

The vector $V := \sum_k \Phi_{k,*}(S_{\tau+1} u_k)$ lies in $\sum_k
\Phi_{k,*}(V_{\mu_{k,*}}) \subseteq E^+_{n-1}|_{\Vlam}$ (Phi-images sit
in $E^+_{n-1}$). By Lemma~\ref{lem:Outer-inj}, $\Outer$ is injective on
$E^+_{n-1}|_{\Vlam}$. From $\Outer V = 0$ we conclude $V = 0$, i.e.,
\[
   \sum_{k \in K_{\hatl}} \Phi_{k,*}(S_{\tau+1} u_k) \;=\; 0.
\]

The $\Phi_{k,*}$-images have pairwise disjoint seminormal supports
(Theorem~\ref{thm:pillar1-red}). Hence each summand vanishes:
$\Phi_{k,*}(S_{\tau+1} u_k) = 0$. Each $\Phi_{k,*}$ is injective, so
$S_{\tau+1} u_k = 0$ for each $k$.

Now apply IH (or, for the base case $\hatl = (2, 2)$,
Lemma~\ref{lem:hook-sharp}) to each $\mu_{k,*}$: by
threshold-matching, $\tau(\hat\mu_{k,*}) = \tau(\hatl)$ for every
contributor (Pillar~1, Lemma~1). So the IH at index $\mu_{k,*}$ says
$S_{\tau+1}\!\mid_{B^{(\mu_{k,*})}}$ is injective, hence $u_k = 0$ for
each $k$.

Therefore $w = \Outer\sum_k \Phi_{k,*}(0) = 0$.\qed
\end{proof}

\begin{remark}[Why threshold-matching is needed here]\label{rem:tau-match}
The IH on $\mu_{k,*}$ gives injectivity of $S_{\tau(\hat\mu_{k,*})+1}$
on $B^{(\mu_{k,*})}$. The threshold-matching identity
(Pillar~1, Lemma~1) guarantees $\tau(\hat\mu_{k,*}) = \tau(\hatl)$ for
every contributor (both regular and shrinking cases). So
$S_{\tau+1}$ is the operator the IH controls. Without
threshold-matching the induction would not close.
\end{remark}

\section{The hook base layer}\label{sec:hook}

\begin{lemma}[Hook sharpness]\label{lem:hook-sharp}
Let $\mu = (k_h, 1^{m_h})$ with $k_h \ge 2$ and $m_h \ge k_h$. Then
$S_{m_h - k_h + 2}$ is injective on $B^{(\mu)}$.
\end{lemma}

\noindent
This is the hook analogue of Theorem~\ref{thm:main} with the
single-row threshold $\tau^{\mathrm{hook}}_0 = m_h - k_h + 1$,
$\tau^{\mathrm{hook}} + 1 = m_h - k_h + 2$.

\begin{remark}\label{rem:hook-tau-match}
Continuing the unfloored convention $\tau_0(\hatm) = q + 2\elh(\hatm)
- 1 - |\hatm|$ to single-row $\hatm = (k_h)$ with $q = m_h$ gives
$\tau_0 = m_h + 2 - 1 - k_h = m_h - k_h + 1$, matching the hook
threshold above. Under the hook-recursion $\mu \to \mu_{1,*} = (k_h-1,
1^{m_h - 1})$, $\tau^{\mathrm{hook}}_0$ is preserved exactly:
$(m_h-1) - (k_h-1) + 1 = m_h - k_h + 1$.
\end{remark}

\begin{proof}[Proof of Lemma~\ref{lem:hook-sharp}]
Induction on $k_h$.

\medskip\textbf{Base case ($k_h = 2$).}
We have $\mu = (2, 1^{m_h})$ with $m_h \ge 2$, $|\mu| = m_h + 2$,
$\ell(\mu) = m_h + 1$, and threshold $m_h - 2 + 2 = m_h$. We must show
$S_{m_h}$ is injective on $B^{(\mu)}$.

The chain $\Om^{(\mu)} = \Rp_{m_h+2}^{(\mu)}$ has a single factor of
length $m_h + 1$, and by Phase~A
\[
   B^{(\mu)} \;=\; \Rp_{m_h+2}^{(\mu)}\, V_\mu
   \;=\; E^+_{m_h+1}\!\mid_{V_\mu}.
\]
The removable cells of $\mu$ are $(1, 2)$ and $(m_h+1, 1)$; the unique
non-column-adjacent cell pair for $\{m_h+1, m_h+2\}$ is the pair
$\{(1, 2), (m_h+1, 1)\}$. Hence $\dim E^+_{m_h+1}|_{V_\mu} = 1$,
spanned by
\[
   w_0 \;:=\; v_{T^+} + \beta\, v_{T^-},
\]
where $T^+$ places $m_h+1$ at $(1, 2)$ and $m_h+2$ at $(m_h+1, 1)$;
$T^-$ swaps the two; and $\beta \in \mathbb{Q}(q)^\times$ is the
Hoefsmit $2$-block coefficient at axial distance $1 - (1 - m_h) =
m_h$.

In $T^+$: entry $m_h$ sits at $(m_h, 1)$, entry $m_h+1$ at $(1, 2)$.
These cells lie in different rows and different columns, so the
$T_{m_h}$-block at $T^+$ is a non-degenerate $2 \times 2$ Hoefsmit
block. In particular, $S_{m_h} v_{T^+} = (a+1) v_{T^+} + b\,
v_{s_{m_h} T^+}$ with $b \ne 0$ (off-diagonal Hoefsmit entry at axial
distance $1 - (1 - m_h) = m_h \ne \pm 1$ is nonzero in $\mathbb{Q}(q)$).

In $T^-$: entry $m_h$ at $(m_h, 1)$, entry $m_h+1$ at $(m_h+1, 1)$ —
same column. So $S_{m_h} v_{T^-} = 0$.

Combining:
\[
   S_{m_h} w_0
   \;=\; (a+1)\, v_{T^+} + b\, v_{s_{m_h} T^+}.
\]
The vector $v_{s_{m_h} T^+}$ is a basis vector outside $\supp(B^{(\mu)})
= \{T^+, T^-\}$ (the SYT $s_{m_h} T^+$ has $m_h$ at $(1, 2)$, so its
prefix length is $m_h - 1 < m_h$; it does not appear in $\supp(B^{(\mu)})$,
in fact it does not appear in the support of $E^+_{m_h+1}|_{V_\mu}$ at
all, but as a basis vector in $V_\mu$). Therefore $S_{m_h} w_0$ has a
nonzero $v_{s_{m_h} T^+}$-coefficient $b \ne 0$, so $S_{m_h} w_0
\ne 0$. Since $B^{(\mu)}$ is the line through $w_0$, this proves
injectivity of $S_{m_h}\!\mid_{B^{(\mu)}}$.

\medskip\textbf{Inductive step ($k_h \ge 3$).}
We invoke the hook analogue of the Pillar~1 reduction theorem,
which we record as Lemma~\ref{lem:hook-reduction} below.

By Lemma~\ref{lem:hook-reduction},
\[
   B^{(\mu)} \;=\; \Outer^{(\mu)} \cdot
   \Phi^{(\mu)}_{1,*}\bigl(B^{(\mu_{1,*})}\bigr),
\]
where $\Outer^{(\mu)} = \prod_{j=m_h+1}^{m_h+k_h-2}(T_j+1)$,
$\mu_{1,*} = (k_h-1, 1^{m_h-1})$, and $\Phi^{(\mu)}_{1,*}$ is the
contributor embedding.

Suppose $S_{m_h-k_h+2} w = 0$ for some $w \in B^{(\mu)}$. Write
$w = \Outer^{(\mu)} \Phi^{(\mu)}_{1,*}(u_1)$ uniquely with $u_1 \in
B^{(\mu_{1,*})}$. We must show $u_1 = 0$.

\emph{Geometric inequalities for the hook.} With
$\ell^{(\mu)} = m_h + 1$ and $\tau^{\mathrm{hook}} + 1 = m_h - k_h + 2$:
\[
   \ell^{(\mu)} - (\tau^{\mathrm{hook}} + 1)
   \;=\; (m_h + 1) - (m_h - k_h + 2) \;=\; k_h - 1 \;\ge\; 2,
\]
and $|\mu| - 1 - (\tau^{\mathrm{hook}} + 1) = (k_h + m_h - 1) -
(m_h - k_h + 2) = 2k_h - 3 \ge 3$. Hence:
\begin{itemize}[topsep=0.2em,itemsep=0.1em]
\item $S_{m_h-k_h+2}$ commutes with every factor of $\Outer^{(\mu)}$
   (each factor index $j \ge m_h + 1 \ge (m_h - k_h + 2) + 2$).
\item $S_{m_h-k_h+2}$ commutes with $T_{|\mu|-1}$, hence preserves
   $E^+_{|\mu|-1}|_{V_\mu}$.
\item $S_{m_h-k_h+2}$ commutes with $\Phi^{(\mu)}_{1,*}$ (since
   $m_h - k_h + 2 \le |\mu| - 3 = k_h + m_h - 3$ iff $2k_h \ge 5$,
   which holds for $k_h \ge 3$).
\end{itemize}

The same argument as in the main theorem proceeds:
\[
   0 \;=\; S_{m_h-k_h+2} w
   \;=\; \Outer^{(\mu)} \Phi^{(\mu)}_{1,*}(S_{m_h-k_h+2} u_1).
\]
The vector $\Phi^{(\mu)}_{1,*}(S_{m_h-k_h+2} u_1)$ lies in
$E^+_{|\mu|-1}|_{V_\mu}$, on which $\Outer^{(\mu)}$ is injective
(Lemma~\ref{lem:Outer-inj} applied at the $\mu$ level — the proof is
identical, since adjacent-injectivity is shape-agnostic). So
$\Phi^{(\mu)}_{1,*}(S_{m_h-k_h+2} u_1) = 0$, and injectivity of
$\Phi^{(\mu)}_{1,*}$ gives $S_{m_h-k_h+2} u_1 = 0$.

The inner threshold matches (Remark~\ref{rem:hook-tau-match}), and the
IH (induction on $k_h$, since $k_h - 1 < k_h$ and $m_h - 1 \ge k_h - 1$
when $m_h \ge k_h \ge 3$) gives $S_{m_h-k_h+2}$ injective on
$B^{(\mu_{1,*})}$. So $u_1 = 0$, hence $w = 0$.\qed
\end{proof}

\begin{lemma}[Hook reduction theorem]\label{lem:hook-reduction}
For $\mu = (k_h, 1^{m_h})$ with $k_h \ge 3$ and $m_h \ge k_h$,
\[
   B^{(\mu)} \;=\; \Outer^{(\mu)} \cdot
   \Phi^{(\mu)}_{1,*}\bigl(B^{(\mu_{1,*})}\bigr),
\]
where $\Outer^{(\mu)} = \prod_{j=m_h+1}^{m_h+k_h-2}(T_j+1)$,
$\mu_{1,*} = (k_h-1, 1^{m_h-1})$, and $\Phi^{(\mu)}_{1,*}$ is the
Phi-embedding for the corner pair $\{(1, k_h), (m_h+1, 1)\}$.
\end{lemma}

\begin{proof}
The proof is identical to the Pillar~1 reduction theorem,
specialised to hooks. We sketch the differences.

\emph{Phi-decomposition of $E^+_{|\mu|-1}|_{V_\mu}$.}
For hook $\mu = (k_h, 1^{m_h})$, the removable corners of $\mu$ are
$(1, k_h)$ and $(m_h + 1, 1)$. The cell pairs for $\{|\mu|-1, |\mu|\}$
are: (i)~the corner pair $\{(1, k_h), (m_h+1, 1)\}$, giving the
contributor $\Phi^{(\mu)}_{1,*}$; (ii)~the same-row pair
$\{(1, k_h - 1), (1, k_h)\}$, giving the vanisher $\Phi^{(\mu)}_h$
when $k_h \ge 3$ (the condition $\hatl_1 \ge 3$ is satisfied; the
condition $\hatl_1 \ge \hatl_2 + 2$ collapses since there is no
$\hatl_2$ for single-row $\hat$ — interpret as automatic).
There are no corner-pair vanishers since $\hatl$ has only one corner.

\emph{Vanisher kill.}
For the same-row vanisher $\Phi^{(\mu)}_h$ at the row-$1$ pair, the
inner shape is $(k_h-2, 1^{m_h})$ if $k_h \ge 4$, or $(1^{m_h+1})$ (a
column) if $k_h = 3$. In either case the hook column-$1$
theorem~\cite[Thm.~4]{pillar1} (or trivially for columns, where
$B = V_\nu \subseteq E^-_1$ since every SYT has $1, 2$ in the same
column) gives $B^{(\nu)} \subseteq E^-_1|_{V_\nu}$. The inner
pull-through provides an extra $\Rp_{\ell}^{(\nu)}$ factor whose
rightmost component is $S_1$, which kills $E^-_1$. So the vanisher
contribution vanishes.

\emph{Contributor.}
For $\mu_{1,*} = (k_h-1, 1^{m_h-1})$, $\ell(\mu_{1,*}) = m_h$, so the
$Q$-product after inner pull-through has factor indices $\ell^{(\mu)}
= m_h + 1$ through $|\mu_{1,*}| = k_h + m_h - 2$, which is exactly
the range $\ell(\mu_{1,*}) + 1, \ldots, |\mu_{1,*}|$ defining
$\Om^{(\mu_{1,*})}$. Hence the contributor yields
$\Phi^{(\mu)}_{1,*}(\Om^{(\mu_{1,*})} V_{\mu_{1,*}}) =
\Phi^{(\mu)}_{1,*}(B^{(\mu_{1,*})})$ after the outer
$\Outer^{(\mu)}$ factors.\qed
\end{proof}

\section{Computational verification}\label{sec:verify}

We verify Theorem~\ref{thm:main} by computing $\dim B^{(\lambda)}$ and
$\dim S_{\tau+1}(B^{(\lambda)})$ for each test shape, using a sparse
Hoefsmit-basis implementation (working mod $P = 100003$ at $Q = 11$).
Sharpness in the strong form is the assertion $\dim S_{\tau+1}(B) =
\dim B$ (full rank of $S_{\tau+1}$ on $B$).

\begin{table}[h]
\centering\renewcommand{\arraystretch}{1.2}
\caption{Verification of strong sharpness on multi-row shapes
$\lambda = (\hatl, 1^q)$ and on hooks $\mu = (k_h, 1^{m_h})$.}
\begin{tabular}{lcccccc}
\toprule
shape & $\elh$ & $n$ & $\tau$ & $\dim V$ & $\dim B$ & $\dim S_{\tau+1}(B)$ \\
\midrule
$(2,2,1^2)$ & $2$ & $6$ & $1$ & $9$ & $1$ & $1$ \\
$(2,2,1^3)$ & $2$ & $7$ & $2$ & $14$ & $1$ & $1$ \\
$(2,2,1^4)$ & $2$ & $8$ & $3$ & $20$ & $1$ & $1$ \\
$(3,2,1^3)$ & $2$ & $8$ & $1$ & $64$ & $2$ & $2$ \\
$(3,2,1^4)$ & $2$ & $9$ & $2$ & $105$ & $2$ & $2$ \\
$(2,2,2,1^2)$ & $3$ & $8$ & $1$ & $28$ & $1$ & $1$ \\
$(2,2,2,1^3)$ & $3$ & $9$ & $2$ & $48$ & $1$ & $1$ \\
$(3,2,2,1^3)$ & $3$ & $10$ & $1$ & $315$ & $3$ & $3$ \\
\midrule
$(2,1^2)$ & --- & $4$ & $1$ & $3$ & $1$ & $1$ \\
$(2,1^3)$ & --- & $5$ & $2$ & $4$ & $1$ & $1$ \\
$(3,1^3)$ & --- & $6$ & $1$ & $10$ & $1$ & $1$ \\
$(3,1^5)$ & --- & $8$ & $3$ & $21$ & $1$ & $1$ \\
$(4,1^4)$ & --- & $8$ & $1$ & $35$ & $1$ & $1$ \\
$(4,1^7)$ & --- & $11$ & $4$ & $120$ & $1$ & $1$ \\
\bottomrule
\end{tabular}
\end{table}

The full rank $\dim S_{\tau+1}(B) = \dim B$ holds in every case
tested.

\section{Discussion}\label{sec:discussion}

\subsection{What this proves}

\paragraph{Sharpness, finally.}
Theorem~\ref{thm:main} closes the structural sharpness gap left open
by the Pillar~1 paper, completing the characterisation of the depth
of the column-$1$ chain image:
\[
   \sup\{j : B^{(\lambda)} \subseteq E^-_j\} \;=\; \tau(\hatl).
\]

\paragraph{Strong form.}
The conclusion is a strict strengthening of the conjectured
$B \not\subseteq E^-_{\tau+1}$: it shows $S_{\tau+1}$ is \emph{injective}
on $B$, i.e., $B \cap E^-_{\tau+1} = 0$. This places $B$ as a complement
to $E^-_{\tau+1}$ inside $\Vlam$, in the sense $B \oplus E^-_{\tau+1}
\hookrightarrow \Vlam$.

\subsection{The role of adjacent-injectivity}

The heart of the proof is Lemma~\ref{lem:Outer-inj}: the outer product
$\Outer = \prod_{j=\ell}^{n-2}(T_j+1)$ is injective on
$E^+_{n-1}|_{\Vlam}$. This statement is module-agnostic — it depends
only on the Hecke quadratic relation and the adjacent-disjointness
$E^-_j \cap E^+_{j+1} = 0$ established in May-19. It is striking that
the same fact that drives r-additivity~\cite{r-add} also drives
sharpness, via the simple observation that $S_{\tau+1}$ preserves
$E^+_{n-1}$.

\subsection{Connection to trace-vanishing}

Combined with Pillar~1, the present theorem says:
\[
   B^{(\lambda)} \;\subseteq\; \bigcap_{j=1}^\tau E^-_j,
   \qquad
   B^{(\lambda)} \,\cap\, E^-_{\tau+1} \;=\; 0.
\]
On $B$, $T_j$ acts as $-1$ for $j \le \tau$ and has a nonzero $+q$
component for $j = \tau + 1$. In particular, the trace
$\mathrm{tr}_{B}(T_1 T_2 \cdots T_{\tau+1}) = (-1)^\tau \cdot
\mathrm{tr}_B(T_{\tau+1})$, where $T_{\tau+1}$ restricted to $B$ has at
least one $+q$ eigenvalue. The exact eigenvalue spectrum of
$T_{\tau+1}|_B$ is a finer invariant.

\subsection{Limitations}

\paragraph{Non-stable regime.}
For $q < |\hatl| - 2\elh + 2$ (the non-stable regime), neither Pillar~1
nor the present theorem applies. The sharpness picture there is
expected to mirror the May-14 non-stable/JM phenomena and is not
addressed here.

\section{Status}\label{sec:status}

\paragraph{Established.}
\begin{itemize}[topsep=0.3em,itemsep=0.1em]
\item Theorem~\ref{thm:main}: $B^{(\lambda)} \cap E^-_{\tau+1} = 0$
   for all $\hatl$ with $\elh \ge 2$, every row $\ge 2$, $q \ge
   |\hatl| - 2\elh + 2$.
\item Lemma~\ref{lem:hook-sharp}: hook sharpness, base layer of the
   recursion.
\item Verification on 14 datapoints (8 multi-row, 6 hook).
\end{itemize}

\paragraph{Open.}
\begin{itemize}[topsep=0.3em,itemsep=0.1em]
\item Eigenvalue spectrum of $T_{\tau+1}|_{B^{(\lambda)}}$.
\item Sharpness in the non-stable regime $q < |\hatl| - 2\elh + 2$.
\item Categorical lift: an analogous statement on the
   triply-graded HOMFLY / Hochschild homology of the Rouquier complex,
   per \texttt{connections/2026-05-16-categorical-home-consolidates.md}.
\end{itemize}

\begin{thebibliography}{9}
\bibitem{pillar1}
   Clio, \emph{Extended sign-kill for general multi-row $\hatl$: the
   meta-theorem at arbitrary $\elh \ge 2$}, 15 May 2026.
   \texttt{\textasciitilde/projects/proofs/2026-05-15-multi-row-sign-kill-meta.tex}.

\bibitem{r-add}
   Clio, \emph{r-additivity and the SYT formula}, 13 May 2026.
   \texttt{\textasciitilde/projects/proofs/2026-05-13-r-additivity-and-syt-formula.tex}.

\bibitem{phaseA}
   Clio, \emph{Phase A endpoint and adjacent-injectivity}, 19 May 2026
   (earlier paper number; cited by Pillar~1).
   \texttt{\textasciitilde/projects/proofs/2026-05-19-phase-a-general.tex}.
\end{thebibliography}

\end{document}
